% !TeX root = ../main.tex
\section{Wnioski}
\label{sec:conclusion}

% ------------------------------------------------------------------------
% Podsumowanie uzyskanych wyników:
%   - wnioski z analizy porównawczej standardów IEEE 802.11a i IEEE 802.11ax
%   - ocena odporności toru przetwarzania na realistyczny clutter kanału TGax Model-B
%   - skuteczność interpolacji nośnych DC i algorytmu Coherent CLEAN
% Ograniczenia zastosowanej metody i implementacji:
%   - aliasing prędkości Dopplera w 802.11ax wynikający z długości symbolu OFDM
%   - stała rozdzielczość odległościowa (7.5 m) przy kanale 20 MHz
%   - czas trwania pojedynczego pakietu (CPI)
% Kierunki dalszych prac:
%   - integracja wielu ramek z kompensacją przerw SIFS/DIFS (Multi-Frame CPI)
%   - wykorzystanie szerszych kanałów 802.11ax (40, 80, 160 MHz) dla wyższej rozdzielczości dR
%   - badania empiryczne z platformami SDR (Software-Defined Radio)
%   - estymacja w oparciu o surowe macierze CSI ze sprzętowych kart sieciowych
% ------------------------------------------------------------------------

\subsection{Podsumowanie uzyskanych wyników}
\label{subsec:conclusion_summary}

\subsection{Ograniczenia zastosowanej metody i implementacji}
\label{subsec:conclusion_limitations}

\subsection{Kierunki dalszych prac}
\label{subsec:conclusion_future_work}